\section{Related Work}
\label{submission:sec:related_work}

\subsection{Parameter-Efficient Fine-Tuning and On-Device Serving}
Parameter-Efficient Fine-Tuning (PEFT) has emerged as the dominant paradigm for adapting large pre-trained foundation models to specialized downstream tasks without incurring the computational and storage overhead of full-parameter fine-tuning. Low-Rank Adaptation (LoRA) \cite{hu2021lora} decomposes weight updates into low-rank matrices, dramatically reducing memory footprints. Serving multiple specialized LoRA adapters simultaneously is highly desirable for personalized on-device applications. Frameworks like S-LoRA \cite{sheng2023slora} and Punica \cite{chen2023punica} address multi-tenant adapter serving by designing custom CUDA kernels (e.g., segmented gather-scatter operations) to batched sequences of heterogeneous inputs. However, these frameworks operate in a multi-pass or context-switched manner, adding scheduling and memory indexing overheads that are often prohibitive for extreme edge processors \cite{lane2015deep, zhou2019edge}.

\subsection{Weight-Space and Activation-Space Model Merging}
An alternative to multi-tenant execution is weight-space and activation-space model merging. Weight merging techniques like Task Arithmetic \cite{ilharco2022editing}, Model Soups \cite{wortsman2022model}, and Ties-Merging \cite{yadav2023ties} linearly combine the parameters of multiple fine-tuned models to construct a single multi-task model. Although computationally zero-cost during inference, weight-space merging suffers from weight interference, where orthogonal or conflicting parameter updates degrade joint performance. To bypass weight-space conflicts, activation-space blending frameworks such as SABLE \cite{mbh2025} and PFSR \cite{pfsr2025} dynamically route and blend activations of specialized adapters. By restricting adapter computations to a single forward pass over a frozen backbone and combining expert activations dynamically, they achieve near-zero interference without scheduling bottlenecks.

\subsection{Coordinate Routing and OOD Task Rejection}
Dynamic model merging architectures rely on rapid, sample-wise task routing. Recent SOTA methods utilize early-stage representations (e.g., Layer 3 CLS token from a Vision Transformer) to bypass the temporal routing paradox \cite{pfsr2025, mbh2025}. In these frameworks, the router projects early features onto pre-computed centroids of registered tasks to map each query into a low-dimensional coordinate space. Out-of-distribution (OOD) task rejection is crucial in these frameworks to prevent arbitrary expert routing under unknown queries. SPS-ZCA \cite{pfsr2025} proposes fitting coordinate-space diagonal GMMs during a calibration phase, rejecting samples that exhibit log-likelihoods below a safe threshold. Despite their conceptual simplicity, the robustness of these coordinate density estimators under non-ideal, noisy deployment conditions has not been systematically investigated.

\subsection{Covariance Estimation and Shrinkage}
The challenge of estimating covariance matrices from small sample sizes is a classical problem in high-dimensional statistics \cite{ledoit2004well}. Under-determined covariance estimates exhibit severe variance underestimation, making probability density estimation highly unstable. Standard regularizers like Ridge regularization (adding a static diagonal ridge $\gamma I$) are simple but lack adaptability, heavily over-regularizing in high-sample regimes and failing to stabilize low-sample estimates. The Ledoit-Wolf shrinkage estimator \cite{ledoit2004well} addresses this by analytically computing the optimal linear combination of the sample covariance matrix and a structured shrinkage target (e.g., a spherical identity matrix). While shrinkage has been applied to linear discriminant analysis and asset pricing, its application to coordinate-space density models for edge model serving represents a novel, mathematically rigorous adaptation to a modern machine learning challenge.
